% Source: https://artofproblemsolving.com/community/c3h2325655_modular_arithmetic_but_its_solvable_without_wolframalpha
% Stable identity: aops-topic-2325655-q01
Factor $2331=9\cdot7\cdot37$, find $N$ modulo $9$, $7$ and $37$, and combine.

\textbf{Modulo $9$.} The digit sum is $71+0+18\cdot9+69\cdot8+9+7+6=71+162+552+22=807$, and $807\equiv8+0+7=15\equiv6$. So $N\equiv6\pmod9$.

\textbf{Three-digit groups.} Since $1000=27\cdot37+1$, we have $1000\equiv1\pmod{37}$; and $1001=7\cdot143$ gives $1000\equiv-1\pmod7$. Split $N$ into $54$ groups of three digits from the left (this matches the grouping from the right, as $162=3\cdot54$): positions $1$--$69$ give $23$ groups $111$; positions $70$--$72$ give $110$; positions $73$--$90$ give $6$ groups $999$; positions $91$--$159$ give $23$ groups $888$; the last group is $976$.

\textbf{Modulo $37$.} $N$ is congruent to the sum of its groups. Since $111=3\cdot37$, $999=27\cdot37$ and $888=24\cdot37$ are multiples of $37$,
\[N\equiv110+976=1086=29\cdot37+13\equiv13\pmod{37}.\]

\textbf{Modulo $7$.} $N$ is congruent to the alternating sum of the groups, with sign $+$ for the last group and alternating towards the left. Modulo $7$: $111\equiv-1$, $110\equiv5$, $999\equiv5$, $888\equiv-1$, $976\equiv3$. The $23$ groups $888$ (groups $31$ to $53$) have signs $-,+,\dots,-$, a net of one $-888\equiv+1$; the $6$ groups $999$ cancel in pairs; the group $110$ (group $24$) has sign $+$; the $23$ groups $111$ (groups $1$ to $23$) have net sign $-$, giving $-111\equiv+1$. So
\[N\equiv1+0+5+1+3=10\equiv3\pmod7.\]

\textbf{Combining.} $x\equiv13\pmod{37}$ and $x\equiv3\pmod7$: $x=13+37t$ with $6+2t\equiv3\pmod7$, so $2t\equiv4$, $t\equiv2$, and $x\equiv87\pmod{259}$. Since $87\equiv6\pmod9$ as well, $N\equiv87\pmod{2331}$. The remainder is $\boxed{87}$.
